\documentclass[12pt]{article}
\usepackage[utf8]{inputenc}
\usepackage[T1]{fontenc}
\usepackage[margin=1in]{geometry}
\usepackage{setspace}
\usepackage{xcolor}
\usepackage[colorlinks=true, urlcolor=blue!60!black, linkcolor=blue!60!black]{hyperref}
\setlength{\parindent}{0em}
\setlength{\parskip}{0.6em}
\renewcommand{\baselinestretch}{1.0}
\begin{document}
\begin{center}
{\LARGE \textit{Brandenburg v. Ohio}, 395 U.S. 444 (1969)\par}
\vspace{0.5em}
{\large Case Brief --- Two-Page Edition\par}
\end{center}

\section{Memory Jogger}
A Klan leader's ``revengeance'' speech was protected: only advocacy directed to and likely to produce imminent lawless action may be punished; \textit{Whitney} overruled (\href{https://supreme.justia.com/cases/federal/us/395/444/}{Justia}).

\section{Detailed Case Facts}
Clarence Brandenburg invited a TV crew to film a Klan rally on a farm attended only by Klan members. Around a burning cross, some armed, he said there ``might have to be some revengeance taken'' if the government continued to ``suppress the white, Caucasian race,'' and announced a march on Congress.

\section{Procedural History}
Convicted under Ohio's Criminal Syndicalism Statute and sentenced to one to ten years, Brandenburg lost in the Ohio courts; the U.S. Supreme Court reversed.

\section{Judicial Votes}
Reversed, 8--0, \emph{per curiam}, with concurrences by Justices Black and Douglas.

\section{Holding}
A State may not proscribe advocacy of force or law violation unless it is directed to inciting or producing imminent lawless action and is likely to do so. Ohio's statute, which punished mere advocacy and assembly, was invalid on its face and as applied.

\section{Analysis of Opinions}
\textbf{Per curiam:} \textit{Whitney} had been ``thoroughly discredited''; quoting \textit{Noto}, the Court separated ``mere abstract teaching'' from ``preparing a group for violent action and steeling it to such action.'' \textbf{Black:} the citation of \textit{Dennis} does not endorse ``clear and present danger.'' \textbf{Douglas:} that test is manipulable; the line is between ``ideas and overt acts.''

\section{Examples: Future Applications}
\textbf{Same-Side:} a pamphlet calling armed revolution a future moral necessity, with no target or date, is protected. \textbf{Opposite-Side:} a bullhorn command to an armed crowd to ``rush the building'' now is unprotected incitement. \textbf{Fence-Sitter:} a livestream telling followers to bring crowbars to a courthouse at noon tomorrow, given \textit{Hess}'s imminence rule.

\section{Critique}
The Court recast \textit{Dennis}'s balancing as an imminence test, never defined ``imminent'' or ``lawless action,'' and adopted a live-crowd paradigm that fits poorly with online radicalization and stochastic incitement.

\section{Key Quotations}
\begin{quote}
``[T]he constitutional guarantees of free speech and free press do not permit a State to forbid or proscribe advocacy of the use of force or of law violation except where such advocacy is directed to inciting or producing imminent lawless action and is likely to incite or produce such action.''
\end{quote}
\noindent --- \emph{Per curiam}, 395 U.S.\ at 447.

\end{document}
